\section{\emph{The Clean Coder}: o profissional}

O \citetitle{livro:martin2011} não é um livro de código; é sobre a
conduta de quem escreve código. Os capítulos mais úteis para a turma:

\begin{description}
\item[Professionalism (cap.~1)] Seções \enquote{Taking Responsibility},
  \enquote{First, Do No Harm} e \enquote{Work Ethic}: responder pelo que
  você entrega, inclusive pelos defeitos, e cuidar da própria formação.
\item[Saying No e Saying Yes (caps.~2 e 3)] \enquote{The Cost of Saying
  Yes} e \enquote{A Language of Commitment}: um \enquote{sim} é um
  compromisso, e negociar o escopo antes custa menos do que prometer e
  não cumprir.
\item[Practicing (cap.~6)] Treinar fora do horário de entrega, como os
  katas da Parte~I.
\item[Estimation e Pressure (caps.~10 e 11)] O capítulo de estimativa
  pergunta \enquote{What Is an Estimate?} e apresenta o PERT, com três
  números em vez de um; o de pressão trata de evitá-la e de lidar com
  ela sem abandonar as disciplinas.
\end{description}

\paragraph{Ligação com a JVM.}
\enquote{Está lento, é o GC} ou \enquote{reescreve com \texttt{StringBuilder}
que fica mais rápido} são promessas sem medição. A postura profissional
que o livro defende pede que a afirmação venha com dado: um perfil do
JFR, uma listagem do \texttt{javap}, um benchmark revisado. Os
capítulos~\ref{cap:jvm-por-dentro} e~\ref{cap:lendo-bytecode} dão as
ferramentas para isso.
